% =====================================================================
%  Projeto 1 - Banco de Dados (Parte 1), Universidade BR
%  Compilar:
%    mkdir -p diagramas/build && dot -Tpdf diagramas/conceitual.dot -o diagramas/build/conceitual.pdf
%    dot -Tpdf diagramas/logico.dot -o diagramas/build/logico.pdf
%    pdflatex projeto-1.tex   (2x)
% =====================================================================
\documentclass[12pt,a4paper]{article}

\usepackage[utf8]{inputenc}
\usepackage[T1]{fontenc}
\usepackage[brazil]{babel}
\usepackage{graphicx}
\usepackage{booktabs}
\usepackage{longtable}
\usepackage{listings}
\usepackage{hyperref}
\usepackage{enumitem}
\usepackage{pdflscape}
\usepackage[margin=2.5cm]{geometry}

\hypersetup{colorlinks=true,linkcolor=blue,urlcolor=blue}
\setlist[itemize]{leftmargin=*,noitemsep}
\setlist[enumerate]{leftmargin=*}
\setlength{\emergencystretch}{3em}

\lstset{
  language=SQL,
  basicstyle=\ttfamily\footnotesize,
  breaklines=true,
  frame=single,
  numbers=left,
  numberstyle=\tiny,
  columns=flexible,
  keepspaces=true,
  showstringspaces=false
}

\title{\textbf{Projeto 1 - Banco de Dados (Parte 1)}\\[0.4em]
\normalsize Universidade BR: modelagem de um sistema de gestão acadêmica}
\author{Mateus Iuri Rocha\\[0.3em]
\normalsize Curso de Tecnologia em Sistemas Inteligentes\\
\normalsize Fatec Capão Bonito}
\date{Entrega: 05/10/2026}

\begin{document}
\maketitle

\begin{abstract}
Este documento apresenta o Projeto de Banco de Dados (Parte 1) da
Universidade BR, contemplando as três fases exigidas pelo enunciado:
modelo conceitual, modelo lógico e modelo físico, além do dicionário de
dados e da verificação da terceira forma normal (3FN). O modelo foi
construído a partir dos treze requisitos do enunciado; o modelo físico é
entregue como script \texttt{.sql} (\texttt{modelo-fisico.sql}) com onze
tabelas e cinco registros em cada uma, executado e conferido no
MySQL/MariaDB.
\end{abstract}

\tableofcontents

% =====================================================================
\section{Introdução}
% =====================================================================

O enunciado pede a construção de um banco de dados para uma universidade,
passando pelas fases:

\begin{enumerate}
  \item Modelo Conceitual (entidades, atributos e relacionamentos);
  \item Modelo Lógico (atributos com tipo de dado, chaves primárias e
        estrangeiras);
  \item Modelo Físico (script \texttt{.sql} com 5 registros em cada
        tabela).
\end{enumerate}

O projeto deve possuir dicionário de dados e estar na 3FN, e todo o
material deve ser documentado em PDF. Este documento é essa entrega, e o
acompanham os arquivos \texttt{modelo-fisico.sql} (script físico),
\texttt{diagramas/conceitual.dot} e \texttt{diagramas/logico.dot}
(fontes dos diagramas, em Graphviz).

% =====================================================================
\section{Regras do enunciado e como foram tratadas}
% =====================================================================

A tabela abaixo mostra cada requisito literal do enunciado e onde ele
aparece no modelo: é a rastreabilidade entre enunciado e entrega.

\begin{longtable}{@{}p{0.46\textwidth}p{0.48\textwidth}@{}}
\toprule
\textbf{Regra do enunciado} & \textbf{Tratamento no modelo} \\
\midrule
\endfirsthead
\toprule
\textbf{Regra do enunciado} & \textbf{Tratamento no modelo} \\
\midrule
\endhead
A universidade BR possui 15 cursos atualmente
& Entidade \texttt{curso}. Foram inseridos 5 cursos de exemplo (mínimo
pedido pelo enunciado); a cardinalidade é independente da carga inicial. \\
O aluno só pode estar matriculado em um curso
& Chave estrangeira \texttt{aluno.codigo\_curso}, cada aluno tem uma
única linha de curso. \\
Os dados do aluno são nome, sobrenome, cpf, rg, filiação, endereço e
telefone
& Todos são atributos de \texttt{aluno} (o endereço foi decomposto em
rua, número, bairro, cidade, uf e cep, ver 1FN na Seção~\ref{sec:3fn}). \\
Os cursos possuem várias disciplinas
& Relação N:N \texttt{curso} $\times$ \texttt{disciplina} materializada
na entidade associativa \texttt{curso\_disciplina}. \\
Disciplinas obrigatórias ou optativas dependendo do curso
& Atributo \texttt{curso\_disciplina.tipo} ('obrigatoria' /
'optativa'): o mesmo par curso--disciplina pode mudar de tipo, por isso
ele fica na associativa e não na disciplina. \\
Pré-requisito para cursar algumas disciplinas
& Auto-relacionamento N:N da disciplina, materializado em
\texttt{disciplina\_pre\_requisito}. \\
O aluno pertence a apenas uma turma
& Entidade \texttt{turma\_ingresso} (a turma de ingresso do curso),
relacionada ao aluno em 1:N, \texttt{aluno.codigo\_turma}. \\
Cada disciplina poderá ter no máximo 40 alunos por turma
& Entidade \texttt{turma} (oferta da disciplina em um semestre) com
\texttt{max\_alunos} $= 40$. Visto Seção~\ref{sec:conceitual}. \\
Os dados das disciplinas: nome, descrição, quantidade de alunos e carga
horária
& Atributos de \texttt{disciplina} (a quantidade de alunos é derivada da
contagem de matrículas, e está documentada como tal no dicionário). \\
A disciplina pertence a departamentos específicos
& Chave estrangeira \texttt{disciplina.codigo\_departamento}, relação
N:1 disciplina--departamento. \\
Alunos podem trancar matricula
& Atributo \texttt{aluno.situacao\_matricula} ('ativa' / 'trancada'). \\
Reprovado no máximo 2 vezes na mesma disciplina
& Restrição de negócio documentada: a contagem de linhas de
\texttt{historico} com \texttt{situacao} = 'reprovado' para o mesmo
aluno e disciplina não pode ultrapassar 2 (ver Seção~\ref{sec:restr}). \\
Histórico escolar com disciplinas cursadas, nota final, frequência e
período
& Entidade \texttt{historico} (aluno $\times$ turma), com
\texttt{nota\_final}, \texttt{frequencia}, \texttt{periodo} e
\texttt{situacao}; a disciplina é alcançada pela turma
(\texttt{turma.codigo\_disciplina}). \\
Professores lecionam no máximo 7 disciplinas ou nenhuma
& Relação N:N \texttt{professor} $\times$ \texttt{disciplina} em
\texttt{professor\_disciplina}; o limite de 7 é restrição documentada. \\
Professores sempre vinculados a um departamento
& \texttt{professor.codigo\_departamento} (N:1, obrigatório). \\
Cada departamento controla um curso específico
& Relação 1:1 \texttt{departamento} $\times$ \texttt{curso}. \\
\bottomrule
\end{longtable}

% =====================================================================
\section{Modelo conceitual}
\label{sec:conceitual}
% =====================================================================

O modelo conceitual (Figura~\ref{fig:conceitual}) possui 7 entidades
fortes, 4 entidades associativas (fruto das relações N:N) e 17
relacionamentos. As entidades fortes são \texttt{curso},
\texttt{departamento}, \texttt{turma\_ingresso}, \texttt{aluno},
\texttt{disciplina}, \texttt{turma} e \texttt{professor}.

Decisões de modelagem que merecem destaque:

\begin{itemize}
  \item \textbf{Duas entidades de turma.} O enunciado diz que ``o aluno
        pertence a apenas uma turma'' e ao mesmo tempo que ``cada
        disciplina poderá ter no máximo 40 alunos por turma''. As duas
        frases se contradizem se turma for a mesma coisa nos dois casos.
        Por isso existem \texttt{turma\_ingresso} (a turma de ingresso
        do curso, em que o aluno se encaixa em 1:N) e \texttt{turma} (a
        oferta de uma disciplina em um semestre, com no máximo 40
        alunos, em que o aluno cursa N:N através do histórico).
  \item \textbf{Pré-requisito como auto-relacionamento.} Uma disciplina
        pode exigir outra disciplina; como o par pode se repetir, a
        relação N:N virou a associativa
        \texttt{disciplina\_pre\_requisito}.
  \item \textbf{Obrigatória/optativa na associativa.} O tipo depende do
        par (curso, disciplina): a mesma disciplina é obrigatória em um
        curso e optativa em outro.
  \item \textbf{Professor--disciplina em N:N.} O enunciado limita o
        professor a 7 disciplinas, mas não diz que a disciplina tem um
        professor só; a relação N:N cobre os dois lados.
\end{itemize}

\begin{landscape}
\begin{figure}[p]
  \centering
  \includegraphics[width=24.7cm,keepaspectratio]%
    {diagramas/build/conceitual.pdf}
  \caption{Modelo conceitual: entidades, atributos e cardinalidades.}
  \label{fig:conceitual}
\end{figure}
\end{landscape}

% =====================================================================
\section{Modelo lógico}
% =====================================================================

A transformação para o modelo lógico (Figura~\ref{fig:logico}) seguiu as
regras clássicas:

\begin{itemize}
  \item cada atributo ganhou tipo de dado (\texttt{integer},
        \texttt{varchar(n)}, \texttt{numeric(4,2)}, \texttt{char(2)});
  \item cada entidade forte recebeu chave primária própria;
  \item as entidades associativas receberam chave primária composta pelas
        chaves das entidades participantes;
  \item cada relação N:1 virou chave estrangeira no lado ``N''
        (\texttt{aluno}, \texttt{disciplina}, \texttt{turma},
        \texttt{professor}, \texttt{turma\_ingresso},
        \texttt{departamento});
  \item a relação 1:1 departamento--curso virou a chave estrangeira
        \texttt{departamento.codigo\_curso}.
\end{itemize}

\begin{landscape}
\begin{figure}[p]
  \centering
  \includegraphics[width=24.7cm,keepaspectratio]%
    {diagramas/build/logico.pdf}
  \caption{Modelo lógico: tipos, chaves primárias e estrangeiras.}
  \label{fig:logico}
\end{figure}
\end{landscape}

% =====================================================================
\section{Modelo físico}
\label{sec:fisico}
% =====================================================================

O script físico está no anexo \texttt{modelo-fisico.sql} e é listado
integralmente a seguir. Ele cria o banco \texttt{ universidade\_br}, as
11 tabelas com suas chaves estrangeiras e carrega \textbf{5 registros em
cada tabela} (55 no total), conforme o enunciado.

Para executar (MySQL/MariaDB):

\begin{lstlisting}[language=bash,frame=none,numbers=none]
sudo mysql universidade_br < modelo-fisico.sql
\end{lstlisting}

Conferência da carga (executada na entrega):

\begin{lstlisting}[language=bash,frame=none,numbers=none]
SELECT 'curso', COUNT(*) FROM universidade_br.curso
UNION ALL SELECT 'departamento', COUNT(*) FROM universidade_br.departamento
UNION ALL SELECT 'turma_ingresso', COUNT(*) FROM universidade_br.turma_ingresso
UNION ALL SELECT 'aluno', COUNT(*) FROM universidade_br.aluno
UNION ALL SELECT 'disciplina', COUNT(*) FROM universidade_br.disciplina
UNION ALL SELECT 'turma', COUNT(*) FROM universidade_br.turma
UNION ALL SELECT 'professor', COUNT(*) FROM universidade_br.professor
UNION ALL SELECT 'curso_disciplina', COUNT(*) FROM universidade_br.curso_disciplina
UNION ALL SELECT 'disciplina_pre_requisito', COUNT(*) FROM universidade_br.disciplina_pre_requisito
UNION ALL SELECT 'professor_disciplina', COUNT(*) FROM universidade_br.professor_disciplina
UNION ALL SELECT 'historico', COUNT(*) FROM universidade_br.historico;
\end{lstlisting}

As 11 linhas retornadas foram \texttt{5} em todas as tabelas.

\subsection{Script completo}

\lstinputlisting{modelo-fisico.sql}

% =====================================================================
\section{Dicionário de dados}
\label{sec:dicionario}
% =====================================================================

\subsection{Primeira parte: visão geral}

\begin{longtable}{@{}ll@{}}
\toprule
\textbf{Campo} & \textbf{Valor} \\
\midrule
\endfirsthead
\toprule
\textbf{Campo} & \textbf{Valor} \\
\midrule
\endhead
Banco de dados & \texttt{universidade\_br} \\
Sistema & Gestão acadêmica da Universidade BR \\
Tabelas & 11 (7 entidades fortes + 4 associativas) \\
Registros por tabela & 5 (55 no total) \\
Autor & Mateus Iuri Rocha \\
Data & 05/10/2026 \\
\bottomrule
\end{longtable}

Relacionamentos existentes:

\begin{longtable}{@{}lll@{}}
\toprule
\textbf{Origem} & \textbf{Destino} & \textbf{Cardinalidade} \\
\midrule
\endfirsthead
\toprule
\textbf{Origem} & \textbf{Destino} & \textbf{Cardinalidade} \\
\midrule
\endhead
departamento & curso & 1:1 \\
turma\_ingresso & curso & N:1 \\
aluno & curso & N:1 \\
aluno & turma\_ingresso & N:1 \\
disciplina & departamento & N:1 \\
turma & disciplina & N:1 \\
professor & departamento & N:1 \\
curso $\times$ disciplina & curso\_disciplina & N:N \\
disciplina $\times$ disciplina & disciplina\_pre\_requisito & N:N \\
professor $\times$ disciplina & professor\_disciplina & N:N \\
aluno $\times$ turma & historico & N:N \\
\bottomrule
\end{longtable}

\subsection{Segunda parte: cada entidade}

\subsubsection{curso}

\begin{longtable}{@{}llll@{}}
\toprule
\textbf{Atributo} & \textbf{Tipo} & \textbf{Chave} & \textbf{Descrição} \\
\midrule
\endfirsthead
\toprule
\textbf{Atributo} & \textbf{Tipo} & \textbf{Chave} & \textbf{Descrição} \\
\midrule
\endhead
codigo & integer & PK & Identificador do curso \\
nome & varchar(60) & --- & Nome do curso \\
turno & varchar(20) & --- & Turno (matutino, vespertino, noturno) \\
\bottomrule
\end{longtable}

\subsubsection{departamento}

\begin{longtable}{@{}llll@{}}
\toprule
\textbf{Atributo} & \textbf{Tipo} & \textbf{Chave} & \textbf{Descrição} \\
\midrule
\endfirsthead
\toprule
\textbf{Atributo} & \textbf{Tipo} & \textbf{Chave} & \textbf{Descrição} \\
\midrule
\endhead
codigo & integer & PK & Identificador do departamento \\
nome & varchar(60) & --- & Nome do departamento \\
codigo\_curso & integer & FK & Curso controlado pelo departamento (1:1) \\
\bottomrule
\end{longtable}

\subsubsection{turma\_ingresso}

\begin{longtable}{@{}llll@{}}
\toprule
\textbf{Atributo} & \textbf{Tipo} & \textbf{Chave} & \textbf{Descrição} \\
\midrule
\endfirsthead
\toprule
\textbf{Atributo} & \textbf{Tipo} & \textbf{Chave} & \textbf{Descrição} \\
\midrule
\endhead
codigo & integer & PK & Identificador da turma de ingresso \\
semestre & varchar(7) & --- & Semestre de ingresso (ex.: 2024.1) \\
codigo\_curso & integer & FK & Curso da turma \\
\bottomrule
\end{longtable}

\subsubsection{aluno}

\begin{longtable}{@{}llll@{}}
\toprule
\textbf{Atributo} & \textbf{Tipo} & \textbf{Chave} & \textbf{Descrição} \\
\midrule
\endfirsthead
\toprule
\textbf{Atributo} & \textbf{Tipo} & \textbf{Chave} & \textbf{Descrição} \\
\midrule
\endhead
cpf & varchar(11) & PK & CPF do aluno (identificador) \\
nome & varchar(60) & --- & Primeiro nome \\
sobrenome & varchar(60) & --- & Sobrenome \\
rg & varchar(20) & --- & Registro Geral \\
filiacao & varchar(120) & --- & Filiação \\
rua & varchar(80) & --- & Rua do endereço \\
numero & varchar(10) & --- & Número do endereço \\
bairro & varchar(60) & --- & Bairro \\
cidade & varchar(60) & --- & Cidade \\
uf & char(2) & --- & Unidade federativa \\
cep & varchar(8) & --- & CEP \\
telefone & varchar(20) & --- & Telefone de contato \\
situacao\_matricula & varchar(20) & --- & 'ativa' ou 'trancada' \\
codigo\_curso & integer & FK & Curso em que o aluno está matriculado \\
codigo\_turma & integer & FK & Turma de ingresso (aluno pertence a 1) \\
\bottomrule
\end{longtable}

\subsubsection{disciplina}

\begin{longtable}{@{}lllp{0.42\textwidth}@{}}
\toprule
\textbf{Atributo} & \textbf{Tipo} & \textbf{Chave} & \textbf{Descrição} \\
\midrule
\endfirsthead
\toprule
\textbf{Atributo} & \textbf{Tipo} & \textbf{Chave} & \textbf{Descrição} \\
\midrule
\endhead
codigo & integer & PK & Identificador da disciplina \\
nome & varchar(60) & --- & Nome da disciplina \\
descricao & varchar(255) & --- & Ementa resumida \\
carga\_horaria & integer & --- & Carga horária em horas \\
quantidade\_alunos & integer & --- & Alunos matriculados (derivado:
contagem de matrículas ativas) \\
codigo\_departamento & integer & FK & Departamento dono da disciplina \\
\bottomrule
\end{longtable}

\subsubsection{turma}

\begin{longtable}{@{}llll@{}}
\toprule
\textbf{Atributo} & \textbf{Tipo} & \textbf{Chave} & \textbf{Descrição} \\
\midrule
\endfirsthead
\toprule
\textbf{Atributo} & \textbf{Tipo} & \textbf{Chave} & \textbf{Descrição} \\
\midrule
\endhead
codigo & integer & PK & Identificador da turma \\
semestre & varchar(7) & --- & Semestre da oferta \\
max\_alunos & integer & --- & Limite de alunos (40) \\
codigo\_disciplina & integer & FK & Disciplina ofertada \\
\bottomrule
\end{longtable}

\subsubsection{professor}

\begin{longtable}{@{}llll@{}}
\toprule
\textbf{Atributo} & \textbf{Tipo} & \textbf{Chave} & \textbf{Descrição} \\
\midrule
\endfirsthead
\toprule
\textbf{Atributo} & \textbf{Tipo} & \textbf{Chave} & \textbf{Descrição} \\
\midrule
\endhead
matricula & integer & PK & Matrícula do professor \\
nome & varchar(60) & --- & Nome completo \\
cpf & varchar(11) & --- & CPF \\
codigo\_departamento & integer & FK & Departamento de vinculação \\
\bottomrule
\end{longtable}

\subsubsection{curso\_disciplina}

\begin{longtable}{@{}llll@{}}
\toprule
\textbf{Atributo} & \textbf{Tipo} & \textbf{Chave} & \textbf{Descrição} \\
\midrule
\endfirsthead
\toprule
\textbf{Atributo} & \textbf{Tipo} & \textbf{Chave} & \textbf{Descrição} \\
\midrule
\endhead
codigo\_curso & integer & PK/FK & Curso da relação \\
codigo\_disciplina & integer & PK/FK & Disciplina da relação \\
tipo & varchar(20) & --- & 'obrigatoria' ou 'optativa' no curso \\
\bottomrule
\end{longtable}

\subsubsection{disciplina\_pre\_requisito}

\begin{longtable}{@{}llll@{}}
\toprule
\textbf{Atributo} & \textbf{Tipo} & \textbf{Chave} & \textbf{Descrição} \\
\midrule
\endfirsthead
\toprule
\textbf{Atributo} & \textbf{Tipo} & \textbf{Chave} & \textbf{Descrição} \\
\midrule
\endhead
codigo\_disciplina & integer & PK/FK & Disciplina que exige o pré-requisito \\
codigo\_pre\_requisito & integer & PK/FK & Disciplina exigida \\
\bottomrule
\end{longtable}

\subsubsection{professor\_disciplina}

\begin{longtable}{@{}llll@{}}
\toprule
\textbf{Atributo} & \textbf{Tipo} & \textbf{Chave} & \textbf{Descrição} \\
\midrule
\endfirsthead
\toprule
\textbf{Atributo} & \textbf{Tipo} & \textbf{Chave} & \textbf{Descrição} \\
\midrule
\endhead
matricula & integer & PK/FK & Professor que leciona \\
codigo\_disciplina & integer & PK/FK & Disciplina lecionada \\
\bottomrule
\end{longtable}

\subsubsection{historico}

\begin{longtable}{@{}llll@{}}
\toprule
\textbf{Atributo} & \textbf{Tipo} & \textbf{Chave} & \textbf{Descrição} \\
\midrule
\endfirsthead
\toprule
\textbf{Atributo} & \textbf{Tipo} & \textbf{Chave} & \textbf{Descrição} \\
\midrule
\endhead
codigo & integer & PK & Identificador do registro \\
cpf\_aluno & varchar(11) & FK & Aluno dono do registro \\
codigo\_turma & integer & FK & Turma cursada \\
nota\_final & numeric(4,2) & --- & Nota final (0 a 10) \\
frequencia & integer & --- & Frequência em percentual \\
periodo & varchar(7) & --- & Período do curso (ex.: 2026.1) \\
situacao & varchar(20) & --- & 'cursando', 'aprovado', 'reprovado' \\
\bottomrule
\end{longtable}

% =====================================================================
\section{Normalização: 1FN, 2FN e 3FN}
\label{sec:3fn}
% =====================================================================

\subsection{Primeira forma normal (1FN)}

Todos os atributos são atômicos: não existem grupos repetidos nem
valores multivalorados. O endereço, que no enunciado aparece como um
campo só, foi decomposto em \texttt{rua}, \texttt{numero},
\texttt{bairro}, \texttt{cidade}, \texttt{uf} e \texttt{cep}; a filiação
e o telefone são um valor por aluno. Cada linha é identificável por uma
chave.

\subsection{Segunda forma normal (2FN)}

Só há chave composta nas entidades associativas, e nelas todo atributo
não-chave depende da chave inteira:

\begin{itemize}
  \item \texttt{curso\_disciplina.tipo} depende do par
        (curso, disciplina): a mesma disciplina é obrigatória em um
        curso e optativa em outro, logo não depende só do curso nem só
        da disciplina;
  \item \texttt{disciplina\_pre\_requisito} e
        \texttt{professor\_disciplina} não têm atributos não-chave;
  \item nas entidades de chave simples, qualquer atributo depende da
        chave inteira por definição.
\end{itemize}

Não há, portanto, dependência parcial.

\subsection{Terceira forma normal (3FN)}

Não há dependência transitiva: nenhum atributo não-chave depende de
outro não-chave:

\begin{itemize}
  \item \texttt{disciplina} guarda o \texttt{codigo\_departamento}, e
        não o nome ou a localização do departamento (que ficariam
        transitivos);
  \item \texttt{aluno} guarda \texttt{codigo\_curso} e
        \texttt{codigo\_turma}, e não o nome do curso;
  \item \texttt{historico} guarda a nota/frequência/periodo, e não
        repete os dados da disciplina, alcançados pela turma;
  \item \texttt{departamento} referencia o curso por chave, mantendo a
        relação 1:1 explícita.
  \item \texttt{disciplina.quantidade\_alunos} é derivado (contagem de
        matrículas); mantivemos o atributo por exigência do enunciado,
        documentado no dicionário como derivado; ele não depende de
        nenhum outro atributo da própria tabela.
\end{itemize}

Todas as 11 tabelas estão na 3FN.

% =====================================================================
\section{Restrições de negócio documentadas}
\label{sec:restr}
% =====================================================================

Algumas regras do enunciado são restrições e não estrutura; foram
deixadas documentadas (podem ser impostas por gatilho/consulta):

\begin{itemize}
  \item no máximo 40 alunos por turma:
        \texttt{SELECT codigo\_turma, COUNT(*) FROM historico
        GROUP BY codigo\_turma HAVING COUNT(*) > 40} deve retornar vazio;
  \item no máximo 7 disciplinas por professor:
        \texttt{SELECT matricula, COUNT(*) FROM professor\_disciplina
        GROUP BY matricula HAVING COUNT(*) > 7};
  \item no máximo 2 reprovações do aluno na mesma disciplina:
        contar linhas de \texttt{historico} com
        \texttt{situacao='reprovado'} agrupadas por
        (cpf\_aluno, disciplina da turma);
  \item professor pode não lecionar nenhuma disciplina: basta não
        existir linha em \texttt{professor\_disciplina};
  \item matrícula trancada: \texttt{aluno.situacao\_matricula}
        \texttt{= 'trancada'}.
\end{itemize}

% =====================================================================
\section{Conclusão}
% =====================================================================

Os três modelos foram entregues e conferidos: o conceitual cobre os 13
requisitos do enunciado, o lógico transforma cada relação em chave
estrangeira ou associativa, e o físico é um script executável que cria
11 tabelas e carrega 5 registros em cada uma. O dicionário de dados
descreve os 52 atributos do banco e a análise mostrou que o modelo está
na 3FN.

\end{document}
